\chapter{Chapter06 — Low Degree Junta Resources}
%%% generated by the blueprint pipeline from FABL.Chapter06.LearningAndTesting.LowDegreeJuntaResources
\section{Overview}

Declarations of the Lean module \texttt{FABL.Chapter06.LearningAndTesting.LowDegreeJuntaResources}.

\section{Declarations}

\begin{definition}[Total random example budget]
\lean{FABL.lowDegreeJuntaTotalRandomExampleBudget}
\label{FABL.lowDegreeJuntaTotalRandomExampleBudget}
\leanok
\uses{FABL.PositiveLearningParameter, FABL.juntaTreeCallCount, FABL.juntaTreePerCallFailure, FABL.lowDegreeJuntaNodeRandomExampleBound}
For natural numbers $n$ and $k$, and a rational failure parameter $\delta$ with
$0<\delta\leq 1/2$, the total random example budget is the number of calls in the
size-$k$ junta tree multiplied by the per-node random example bound for dimension $n$,
size parameter $k$, and the per-call failure parameter assigned from $k$ and $\delta$.
\end{definition}

\begin{definition}[Total work budget]
\lean{FABL.lowDegreeJuntaTotalWorkBudget}
\label{FABL.lowDegreeJuntaTotalWorkBudget}
\leanok
\uses{FABL.PositiveLearningParameter, FABL.juntaTreeCallCount, FABL.juntaTreePerCallFailure, FABL.lowDegreeJuntaNodeUniformWorkBound}
For natural numbers $n$ and $k$, and a rational failure parameter $\delta$ with
$0<\delta\leq 1/2$, the total work budget is the number of calls in the size-$k$ junta
tree multiplied by the per-node uniform work bound for dimension $n$, size parameter
$k$, and the per-call failure parameter assigned from $k$ and $\delta$.
\end{definition}

\begin{theorem}[Total cost of the junta learner]
\lean{FABL.lowDegreeJuntaLearnerProgram_cost_le_totalBudget}
\label{FABL.lowDegreeJuntaLearnerProgram_cost_le_totalBudget}
\leanok
\uses{FABL.BooleanFunction, FABL.DecisionTree, FABL.LearningCost, FABL.LearningProgram.runWithCost, FABL.PositiveLearningParameter, FABL.Sign, FABL.lowDegreeJuntaLearnerProgram, FABL.lowDegreeJuntaLearnerProgram_cost_le, FABL.lowDegreeJuntaTotalRandomExampleBudget, FABL.lowDegreeJuntaTotalWorkBudget}
\difficulty{1}
Let $n,k\in\mathbb{N}$, let $f:\{\pm1\}^n\to\{\pm1\}$ be a Boolean function, and let
$\delta\in\mathbb{Q}$ satisfy $0<\delta\le 1/2$. For any outcome in the support of the
learner's run on $f$, the number of random examples is at most its total random-example
budget, the number of membership queries is zero, and its local work is at most its
total work budget.
\end{theorem}

\begin{definition}[Fourier scan runtime envelope]
\lean{FABL.lowDegreeJuntaFourierScanRuntimeEnvelope}
\label{FABL.lowDegreeJuntaFourierScanRuntimeEnvelope}
\leanok
\uses{FABL.PositiveLearningParameter, FABL.fourierEstimatorFailureBits, FABL.lowDegreeJuntaCoefficientAccuracy, FABL.lowDegreeJuntaCoefficientConfidence, FABL.lowDegreeJuntaCutoff}
Let $n$ and $k$ be natural numbers, let $P$ be a subset of the $n$ coordinates, and let
$\delta$ be a rational failure parameter with $0<\delta\leq 1/2$. Write $c$ for the
low-degree cutoff at $k$, $a$ for the coefficient accuracy at $k$, $b$ for the
failure-bit count at the coefficient confidence assigned to $P,k,\delta$, and $q$ for
that confidence's per-coordinate failure parameter. The Fourier scan runtime envelope is
$((c+1)(n+1)^c)\, (4bq/a^2)\, (c+1)$.
\end{definition}

\begin{definition}[Polynomial runtime envelope]
\lean{FABL.lowDegreeJuntaPolynomialRuntimeEnvelope}
\label{FABL.lowDegreeJuntaPolynomialRuntimeEnvelope}
\leanok
\uses{FABL.PositiveLearningParameter, FABL.fourierEstimatorFailureBits, FABL.lowDegreeJuntaPolynomialDegree, FABL.lowDegreeJuntaQuarterFailure}
For natural numbers $n$ and $k$, and a rational failure parameter $\delta$ with
$0<\delta\leq 1/2$, let $d$ be the low-degree junta polynomial degree at $k$, and let
$b$ be the failure-bit count at the quarter-failure parameter associated to $\delta$.
The polynomial runtime envelope is $3\bigl(2^d(d+b+3)\bigr)^3(n+1)^{3d}$.
\end{definition}

\begin{theorem}[Tree-Scaled Runtime Ledger Bounds]
\lean{FABL.lowDegreeJuntaNodeRuntimeLedger_mul_treeCallCount}
\label{FABL.lowDegreeJuntaNodeRuntimeLedger_mul_treeCallCount}
\leanok
\uses{FABL.JuntaFreeAssignment, FABL.PositiveLearningParameter, FABL.Sign, FABL.juntaRestrictionDimension, FABL.juntaTreeCallCount, FABL.lowDegreeF₂PolynomialLearnerWork, FABL.lowDegreeJuntaCoefficientSampleCount, FABL.lowDegreeJuntaCutoff, FABL.lowDegreeJuntaFourierFamily, FABL.lowDegreeJuntaFourierScanRuntimeEnvelope, FABL.lowDegreeJuntaNodeRuntimeLedger, FABL.lowDegreeJuntaPolynomialDegree, FABL.lowDegreeJuntaPolynomialRuntimeEnvelope, FABL.lowDegreeJuntaPolynomialSampleCount, FABL.lowDegreeJuntaPolynomialSamples, FABL.lowDegreeJuntaRestrictedSampleCount}
\difficulty{4}
Let $n,k\in\mathbb N$, let $P\subseteq[n]$, let $\delta\in\mathbb Q$ satisfy
$0<\delta\le1/2$, and let a batch of the prescribed restricted sample size assign a free
sign configuration and an output sign to each sample index. If $T_0=1$ and
$T_{j+1}=1+2T_j$, then the tree call count $T_k$ times the Fourier scan work is at most
$T_k$ times the Fourier scan runtime envelope. Also, $T_k$ times the polynomial
learner's work on the reindexed polynomial sample batch is at most $T_k$ times the
polynomial runtime envelope.
\end{theorem}

\begin{definition}[Dyadic failure parameter]
\lean{FABL.lowDegreeJuntaDyadicFailure}
\label{FABL.lowDegreeJuntaDyadicFailure}
\leanok
\uses{FABL.PositiveLearningParameter}
For a natural number $m$, the dyadic failure parameter is $2^{-(m+1)}$.
\end{definition}

\begin{theorem}[Dyadic Failure Parameter Value]
\lean{FABL.lowDegreeJuntaDyadicFailure_value}
\label{FABL.lowDegreeJuntaDyadicFailure_value}
\leanok
\uses{FABL.lowDegreeJuntaDyadicFailure}
\difficulty{1}
For every natural number $b$, the rational value of the dyadic failure parameter with
$b$ confidence bits is $2^{-(b+1)}$.
\end{theorem}

\begin{theorem}[Dyadic form of the failure parameter]
\lean{FABL.lowDegreeJuntaLearningFailure_eq_dyadicFailure}
\label{FABL.lowDegreeJuntaLearningFailure_eq_dyadicFailure}
\leanok
\uses{FABL.lowDegreeJuntaDyadicFailure, FABL.lowDegreeJuntaLearningFailure}
\difficulty{3}
The low-degree junta learning failure parameter equals the dyadic failure parameter with
three bits; both are $1/16$.
\end{theorem}

\begin{theorem}[Failure bits for dyadic parameters]
\lean{FABL.fourierEstimatorFailureBits_lowDegreeJuntaDyadicFailure}
\label{FABL.fourierEstimatorFailureBits_lowDegreeJuntaDyadicFailure}
\leanok
\uses{FABL.fourierEstimatorFailureBits, FABL.lowDegreeJuntaDyadicFailure, FABL.lowDegreeJuntaDyadicFailure_value}
\difficulty{4}
For every natural number $b$, take the learning parameter $\delta=1/2^{b+1}$. Its
failure-bit count $\lceil\log_2\lceil 2/\delta\rceil\rceil$ is exactly $b+2$.
\end{theorem}

\begin{definition}[Dyadic per-node failure parameter]
\lean{FABL.lowDegreeJuntaDyadicNodeFailure}
\label{FABL.lowDegreeJuntaDyadicNodeFailure}
\leanok
\uses{FABL.PositiveLearningParameter, FABL.juntaTreePerCallFailure, FABL.lowDegreeJuntaDyadicFailure}
For natural numbers $k$ and $m$, the dyadic per-node failure parameter is the per-call
failure allocation for a size-$k$ junta tree when the overall dyadic failure parameter
is $2^{-(m+1)}$.
\end{definition}

\begin{theorem}[Dyadic Node Failure Parameter Value]
\lean{FABL.lowDegreeJuntaDyadicNodeFailure_value}
\label{FABL.lowDegreeJuntaDyadicNodeFailure_value}
\leanok
\uses{FABL.juntaTreeCallCount, FABL.juntaTreeCallCount_pos, FABL.juntaTreePerCallFailure_value, FABL.lowDegreeJuntaDyadicFailure_value, FABL.lowDegreeJuntaDyadicNodeFailure}
\difficulty{4}
Let $k$ and $b$ be natural numbers, and let $T_0=1$ and $T_{j+1}=1+2T_j$. The rational
value of the dyadic node failure parameter is $1/(T_k2^{b+1})$.
\end{theorem}

\begin{theorem}[Failure bits for reciprocal parameters]
\lean{FABL.fourierEstimatorFailureBits_eq_clog_two_mul_denominator}
\label{FABL.fourierEstimatorFailureBits_eq_clog_two_mul_denominator}
\leanok
\uses{FABL.PositiveLearningParameter, FABL.fourierEstimatorFailureBits}
\difficulty{3}
Let $\delta$ be a rational number with $0<\delta\leq 1/2$, and let $d$ be a positive
natural number. If $\delta=1/d$, then the failure-bit count $\lceil\log_2\lceil
2/\delta\rceil\rceil$ equals $\lceil\log_2(2d)\rceil$.
\end{theorem}

\begin{theorem}[Binary logarithm of a product]
\lean{FABL.lowDegreeJuntaBinaryClog_mul_le}
\label{FABL.lowDegreeJuntaBinaryClog_mul_le}
\leanok
\difficulty{3}
For all natural numbers $a$ and $b$, the base-two ceiling logarithm satisfies
$\lceil\log_2(ab)\rceil\leq \lceil\log_2 a\rceil+\lceil\log_2 b\rceil$, with the
convention for the ceiling logarithm at zero given by its standard definition.
\end{theorem}

\begin{theorem}[Size bound for a low-degree Fourier family]
\lean{FABL.clog_card_lowDegreeJuntaFourierFamily_le}
\label{FABL.clog_card_lowDegreeJuntaFourierFamily_le}
\leanok
\uses{FABL.card_lowDegreeJuntaFourierFamily_le, FABL.lowDegreeJuntaCutoff, FABL.lowDegreeJuntaFourierFamily}
\difficulty{4}
Let $n,k$ be natural numbers and let $P$ be a subset of the $n$ coordinates. Form the
Fourier family on the coordinates outside $P$, using degrees at most $\lfloor
3k/4\rfloor$. The base-two logarithm rounded up of the number of sets in this family is
at most $(\lfloor 3k/4\rfloor+1)(n+2)$.
\end{theorem}

\begin{definition}[Linear confidence coefficient]
\lean{FABL.lowDegreeJuntaLinearConfidenceCoefficient}
\label{FABL.lowDegreeJuntaLinearConfidenceCoefficient}
\leanok
\uses{FABL.juntaTreeCallCount, FABL.lowDegreeJuntaCutoff}
For each natural number $k$, let $T_k$ be the junta tree call count and $D_k$ the low
degree junta cutoff. The linear confidence coefficient is the natural number $\lceil
\log_2 T_k\rceil + 2(D_k+1)+5$, where the first term denotes the least natural number
$m$ such that $T_k\le 2^m$.
\end{definition}

\begin{theorem}[Half-node failure-bit bound]
\lean{FABL.lowDegreeJuntaHalfNodeFailureBits_le}
\label{FABL.lowDegreeJuntaHalfNodeFailureBits_le}
\leanok
\uses{FABL.fourierEstimatorFailureBits, FABL.fourierEstimatorFailureBits_eq_clog_two_mul_denominator, FABL.juntaTreeCallCount, FABL.juntaTreeCallCount_pos, FABL.lowDegreeJuntaBinaryClog_mul_le, FABL.lowDegreeJuntaDyadicNodeFailure, FABL.lowDegreeJuntaDyadicNodeFailure_value, FABL.lowDegreeJuntaHalfFailure, FABL.lowDegreeJuntaHalfFailure_value, FABL.lowDegreeJuntaLinearConfidenceCoefficient}
\difficulty{5}
For natural numbers $k$ and $b$, let $T_0=1$ and $T_{j+1}=1+2T_j$, and let
$\delta=1/2^{b+1}$. The number of failure bits associated with the parameter
$\delta/(2T_k)$ is at most $\bigl(\lceil\log_2 T_k\rceil+2(\lfloor
3k/4\rfloor+1)+5\bigr)(b+1)$, where the number of failure bits for a positive parameter
$q$ is $\lceil\log_2\lceil 2/q\rceil\rceil$.
\end{theorem}

\begin{theorem}[Quartered Node Failure Bits Bound]
\lean{FABL.lowDegreeJuntaQuarterNodeFailureBits_le}
\label{FABL.lowDegreeJuntaQuarterNodeFailureBits_le}
\leanok
\uses{FABL.fourierEstimatorFailureBits, FABL.fourierEstimatorFailureBits_eq_clog_two_mul_denominator, FABL.juntaTreeCallCount, FABL.juntaTreeCallCount_pos, FABL.lowDegreeJuntaBinaryClog_mul_le, FABL.lowDegreeJuntaDyadicNodeFailure, FABL.lowDegreeJuntaDyadicNodeFailure_value, FABL.lowDegreeJuntaLinearConfidenceCoefficient, FABL.lowDegreeJuntaQuarterFailure, FABL.lowDegreeJuntaQuarterFailure_value}
\difficulty{5}
For all $k,b\in\mathbb{N}$, let $T_0=1$ and $T_{j+1}=1+2T_j$, and let
$\delta=1/(4T_k2^{b+1})$. The failure-bit count $\left\lceil\log_2\left\lceil
2/\delta\right\rceil\right\rceil$ is at most
$\bigl(\lceil\log_2T_k\rceil+2(\lfloor3k/4\rfloor+1)+5\bigr)(b+1)$.
\end{theorem}

\begin{theorem}[Linear bound on coefficient confidence bits]
\lean{FABL.lowDegreeJuntaCoefficientFailureBits_le_linearConfidence}
\label{FABL.lowDegreeJuntaCoefficientFailureBits_le_linearConfidence}
\leanok
\uses{FABL.clog_card_lowDegreeJuntaFourierFamily_le, FABL.fourierEstimatorFailureBits, FABL.fourierEstimatorFailureBits_eq_clog_two_mul_denominator, FABL.juntaTreeCallCount, FABL.juntaTreeCallCount_pos, FABL.lowDegreeJuntaBinaryClog_mul_le, FABL.lowDegreeJuntaCoefficientConfidence, FABL.lowDegreeJuntaCoefficientConfidence_value, FABL.lowDegreeJuntaCutoff, FABL.lowDegreeJuntaDyadicNodeFailure, FABL.lowDegreeJuntaDyadicNodeFailure_value, FABL.lowDegreeJuntaFourierFamily, FABL.lowDegreeJuntaFourierFamily_nonempty, FABL.lowDegreeJuntaLinearConfidenceCoefficient}
\difficulty{6}
Let $n,k,b$ be natural numbers, and let $P\subseteq[n]$. Define $T_0=1$ and
$T_{j+1}=1+2T_j$, and put $r=\lfloor3k/4\rfloor$ and $L=\lceil\log_2
T_k\rceil+2(r+1)+5$. For a failure parameter $\delta=1/(2^{b+1}T_k)$ and coefficient
confidence $\delta/(4|\mathcal F|)$, where $\mathcal F$ is the family of subsets of the
coordinates outside $P$ of size at most $r$, the number of failure bits is at most
$L(n+1)(b+1)$.
\end{theorem}

\begin{theorem}[Coefficient sample count bound]
\lean{FABL.lowDegreeJuntaCoefficientSampleCount_le_failureBits}
\label{FABL.lowDegreeJuntaCoefficientSampleCount_le_failureBits}
\leanok
\uses{FABL.PositiveLearningParameter, FABL.fourierEstimatorFailureBits, FABL.fourierEstimatorSampleCount_cast_le, FABL.lowDegreeJuntaCoefficientAccuracy, FABL.lowDegreeJuntaCoefficientConfidence, FABL.lowDegreeJuntaCoefficientSampleCount}
\difficulty{3}
Let $n$ be a natural number, let $P\subseteq[n]$, let $k$ be a natural number, and let
$\delta\in\mathbb{Q}$ satisfy $0<\delta\leq1/2$. Let $\mathcal F$ be the family of
subsets of the coordinates outside $P$ of size at most $\lfloor3k/4\rfloor$, and set
$\gamma=\delta/(4|\mathcal F|)$. The coefficient sample count for accuracy
$1/(3\cdot2^k)$ and confidence $\gamma$ is at most $36\cdot4^k$ times the number of
failure bits associated with $\gamma$.
\end{theorem}

\begin{definition}[Node base resource exponent]
\lean{FABL.lowDegreeJuntaNodeBaseResourceExponent}
\label{FABL.lowDegreeJuntaNodeBaseResourceExponent}
\leanok
\uses{FABL.lowDegreeJuntaCutoff}
For each natural number $k$, let $D_k$ be the low degree junta cutoff. The node base
resource exponent is $D_k+3$.
\end{definition}

\begin{definition}[Node work resource exponent]
\lean{FABL.lowDegreeJuntaNodeWorkResourceExponent}
\label{FABL.lowDegreeJuntaNodeWorkResourceExponent}
\leanok
\uses{FABL.lowDegreeJuntaCutoff}
For each natural number $k$, let $D_k$ be the low degree junta cutoff. The node work
resource exponent is $D_k+4$.
\end{definition}

\begin{definition}[Fourier batch linear coefficient]
\lean{FABL.lowDegreeJuntaFourierBatchLinearCoefficient}
\label{FABL.lowDegreeJuntaFourierBatchLinearCoefficient}
\leanok
\uses{FABL.lowDegreeJuntaCutoff, FABL.lowDegreeJuntaLinearConfidenceCoefficient}
For each natural number $k$, let $D_k$ be the low degree junta cutoff and let $C_k$ be
the linear confidence coefficient. The Fourier batch linear coefficient is
$(D_k+1)(36\cdot 4^k\cdot C_k)$.
\end{definition}

\begin{definition}[Polynomial sample linear coefficient]
\lean{FABL.lowDegreeJuntaPolynomialSampleLinearCoefficient}
\label{FABL.lowDegreeJuntaPolynomialSampleLinearCoefficient}
\leanok
\uses{FABL.lowDegreeJuntaLinearConfidenceCoefficient, FABL.lowDegreeJuntaPolynomialDegree}
For each natural number $k$, let $P_k$ be the low degree junta polynomial degree and let
$C_k$ be the linear confidence coefficient. The polynomial sample linear coefficient is
$2^{P_k}(P_k+1+C_k)$.
\end{definition}

\begin{definition}[Restricted sample linear coefficient]
\lean{FABL.lowDegreeJuntaRestrictedSampleLinearCoefficient}
\label{FABL.lowDegreeJuntaRestrictedSampleLinearCoefficient}
\leanok
\uses{FABL.lowDegreeJuntaFourierBatchLinearCoefficient, FABL.lowDegreeJuntaPolynomialSampleLinearCoefficient}
For each natural number $k$, let $F_k$ be the Fourier batch linear coefficient and $S_k$
the polynomial sample linear coefficient. The restricted sample linear coefficient is
$F_k+S_k$.
\end{definition}

\begin{definition}[Polynomial work linear coefficient]
\lean{FABL.lowDegreeJuntaPolynomialWorkLinearCoefficient}
\label{FABL.lowDegreeJuntaPolynomialWorkLinearCoefficient}
\leanok
\uses{FABL.lowDegreeJuntaPolynomialDegree, FABL.lowDegreeJuntaPolynomialSampleLinearCoefficient}
For each natural number $k$, let $P_k$ be the low degree junta polynomial degree and let
$S_k$ be the polynomial sample linear coefficient. The polynomial work linear
coefficient is $S_k(P_k+2)+(P_k+1)S_k(P_k+3)+(P_k+1)(P_k+2)$.
\end{definition}

\begin{definition}[Node analysis linear coefficient]
\lean{FABL.lowDegreeJuntaNodeAnalysisLinearCoefficient}
\label{FABL.lowDegreeJuntaNodeAnalysisLinearCoefficient}
\leanok
\uses{FABL.lowDegreeJuntaCutoff, FABL.lowDegreeJuntaFourierBatchLinearCoefficient, FABL.lowDegreeJuntaPolynomialDegree, FABL.lowDegreeJuntaPolynomialWorkLinearCoefficient}
For each natural number $k$, let $F_k$ be the Fourier batch linear coefficient, $D_k$
the low degree junta cutoff, $W_k$ the polynomial work linear coefficient, and $P_k$ the
low degree junta polynomial degree. The node analysis linear coefficient is
$F_k(D_k+1)+W_k+(D_k+1)+(P_k+1)$.
\end{definition}

\begin{definition}[Random example linear coefficient]
\lean{FABL.lowDegreeJuntaNodeRandomExampleLinearCoefficient}
\label{FABL.lowDegreeJuntaNodeRandomExampleLinearCoefficient}
\leanok
\uses{FABL.lowDegreeJuntaLinearConfidenceCoefficient, FABL.lowDegreeJuntaRestrictedSampleLinearCoefficient}
For each natural number $k$, define the random-example linear coefficient to be
$2^{k+1}$ times the restricted-sample linear coefficient at $k$, plus $16\cdot 4^k$
times the linear confidence coefficient at $k$.
\end{definition}

\begin{definition}[Node work linear coefficient]
\lean{FABL.lowDegreeJuntaNodeWorkLinearCoefficient}
\label{FABL.lowDegreeJuntaNodeWorkLinearCoefficient}
\leanok
\uses{FABL.lowDegreeJuntaNodeAnalysisLinearCoefficient, FABL.lowDegreeJuntaNodeRandomExampleLinearCoefficient, FABL.lowDegreeJuntaRestrictedSampleLinearCoefficient}
For each natural number $k$, define the node-work linear coefficient to be twice the
node random-example linear coefficient at $k$, plus the restricted-sample linear
coefficient and the node-analysis linear coefficient at $k$.
\end{definition}

\begin{theorem}[Fourier Batch Size Bound]
\lean{FABL.lowDegreeJuntaFourierBatchSize_le_linearConfidence}
\label{FABL.lowDegreeJuntaFourierBatchSize_le_linearConfidence}
\leanok
\uses{FABL.card_lowDegreeJuntaFourierFamily_le, FABL.lowDegreeJuntaCoefficientFailureBits_le_linearConfidence, FABL.lowDegreeJuntaCoefficientSampleCount, FABL.lowDegreeJuntaCoefficientSampleCount_le_failureBits, FABL.lowDegreeJuntaCutoff, FABL.lowDegreeJuntaDyadicNodeFailure, FABL.lowDegreeJuntaFourierBatchLinearCoefficient, FABL.lowDegreeJuntaFourierFamily, FABL.lowDegreeJuntaLinearConfidenceCoefficient}
\difficulty{3}
Let $n,k,b$ be natural numbers, let $P\subseteq\{0,\ldots,n-1\}$, and let $\delta$ be
rational with $0<\delta\le 1/2$. Write $d=n-|P|$, $r=\lfloor 3k/4\rfloor$, and let
$\mathcal F$ be the family of subsets of a $d$-element set having size at most $r$. Put
$T_0=1$, $T_{j+1}=1+2T_j$, $c=\lceil\log_2\lceil 2/\delta\rceil\rceil$, $L=c+2(r+1)+5$,
and $A=(r+1)(36\cdot4^kL)$. Define $\alpha=1/(3\cdot2^k)$, $\eta=\delta/(4|\mathcal
F|T_k)$, and let $S=\left\lceil
2\,\lceil\log_2\lceil2/\eta\rceil\rceil/\alpha^2\right\rceil$. Then $|\mathcal F|S\le
A(n+1)^{r+1}(b+1)$.
\end{theorem}

\begin{theorem}[Polynomial Sample Count Bound]
\lean{FABL.lowDegreeJuntaPolynomialSampleCount_le_linearConfidence}
\label{FABL.lowDegreeJuntaPolynomialSampleCount_le_linearConfidence}
\leanok
\uses{FABL.fourierEstimatorFailureBits, FABL.juntaRestrictionDimension, FABL.juntaRestrictionDimension_le, FABL.lowDegreeF₂MonomialCount, FABL.lowDegreeF₂MonomialCount_le, FABL.lowDegreeJuntaDyadicNodeFailure, FABL.lowDegreeJuntaLinearConfidenceCoefficient, FABL.lowDegreeJuntaPolynomialDegree, FABL.lowDegreeJuntaPolynomialSampleCount, FABL.lowDegreeJuntaPolynomialSampleLinearCoefficient, FABL.lowDegreeJuntaQuarterFailure, FABL.lowDegreeJuntaQuarterNodeFailureBits_le}
\difficulty{4}
For every $n\in\mathbb{N}$, every subset $P$ of the $n$ coordinates, and all
$k,b\in\mathbb{N}$, set $\ell=k-\lfloor 3k/4\rfloor$ and $T_0=1$, $T_{j+1}=1+2T_j$. Let
$\delta$ be the per-node failure parameter obtained by dividing $2^{-(b+1)}$ by $T_k$.
The polynomial learner's sample count at failure parameter $\delta$ is at most
$2^\ell\bigl(\ell+1+\lceil\log_2
T_k\rceil+2(\lfloor3k/4\rfloor+1)+5\bigr)(n+1)^\ell(b+1)$.
\end{theorem}

\begin{theorem}[Restricted Sample Count Bound]
\lean{FABL.lowDegreeJuntaRestrictedSampleCount_le_linearConfidence}
\label{FABL.lowDegreeJuntaRestrictedSampleCount_le_linearConfidence}
\leanok
\uses{FABL.lowDegreeJuntaCoefficientSampleCount, FABL.lowDegreeJuntaCutoff, FABL.lowDegreeJuntaDyadicNodeFailure, FABL.lowDegreeJuntaFourierBatchLinearCoefficient, FABL.lowDegreeJuntaFourierBatchSize_le_linearConfidence, FABL.lowDegreeJuntaFourierFamily, FABL.lowDegreeJuntaNodeBaseResourceExponent, FABL.lowDegreeJuntaPolynomialDegree, FABL.lowDegreeJuntaPolynomialSampleCount, FABL.lowDegreeJuntaPolynomialSampleCount_le_linearConfidence, FABL.lowDegreeJuntaPolynomialSampleLinearCoefficient, FABL.lowDegreeJuntaRestrictedSampleCount, FABL.lowDegreeJuntaRestrictedSampleLinearCoefficient, FABL.lowDegreeJuntaRuntimeExponentLedger}
\difficulty{4}
For every $n\in\mathbb{N}$, every subset $P$ of the $n$ coordinates, and all
$k,b\in\mathbb{N}$, set $q=\lfloor3k/4\rfloor$, $\ell=k-q$, and $T_0=1$,
$T_{j+1}=1+2T_j$. Put
$A=(q+1)36\cdot4^k\bigl(\lceil\log_2T_k\rceil+2(q+1)+5\bigr)+2^\ell\bigl(\ell+1+\lceil\log_2T_k\rceil+2(q+1)+5\bigr)$.
The restricted sample count, at failure parameter $2^{-(b+1)}/T_k$, is at most
$A(n+1)^{q+3}(b+1)$.
\end{theorem}

\begin{theorem}[Polynomial Learner Work Bound]
\lean{FABL.lowDegreeJuntaPolynomialWorkBound_le_linearConfidence}
\label{FABL.lowDegreeJuntaPolynomialWorkBound_le_linearConfidence}
\leanok
\uses{FABL.juntaRestrictionDimension, FABL.juntaRestrictionDimension_le, FABL.lowDegreeF₂MonomialCount, FABL.lowDegreeF₂MonomialCount_le, FABL.lowDegreeJuntaDyadicNodeFailure, FABL.lowDegreeJuntaPolynomialDegree, FABL.lowDegreeJuntaPolynomialSampleCount, FABL.lowDegreeJuntaPolynomialSampleCount_le_linearConfidence, FABL.lowDegreeJuntaPolynomialSampleLinearCoefficient, FABL.lowDegreeJuntaPolynomialWorkBound, FABL.lowDegreeJuntaPolynomialWorkLinearCoefficient}
\difficulty{5}
For every $n\in\mathbb{N}$, every subset $P$ of the $n$ coordinates, and all
$k,b\in\mathbb{N}$, set $\ell=k-\lfloor3k/4\rfloor$ and $T_0=1$, $T_{j+1}=1+2T_j$.
Define $C=2^\ell(\ell+1+\lceil\log_2T_k\rceil+2(\lfloor3k/4\rfloor+1)+5)$,
$D=\binom{r}{\le\ell}$ where $r$ is the number of coordinates outside $P$, and $m$ to be
the polynomial sample count at failure parameter $2^{-(b+1)}/T_k$. Then the work bound
$m(D+1)+Dm(D+2)+D(D+1)$ is at most
$\bigl(C(\ell+2)+(\ell+1)C(\ell+3)+(\ell+1)(\ell+2)\bigr)(n+1)^{3\ell}(b+1)$. Here
$\binom{r}{\le\ell}$ denotes the number of subsets of an $r$-element set of size at most
$\ell$.
\end{theorem}

\begin{theorem}[Node analysis work bound]
\lean{FABL.lowDegreeJuntaNodeAnalysisWorkBound_le_linearConfidence}
\label{FABL.lowDegreeJuntaNodeAnalysisWorkBound_le_linearConfidence}
\leanok
\uses{FABL.card_lowDegreeJuntaFourierFamily_le, FABL.juntaRestrictionDimension, FABL.juntaRestrictionDimension_le, FABL.lowDegreeF₂MonomialCount, FABL.lowDegreeF₂MonomialCount_le, FABL.lowDegreeJuntaCoefficientSampleCount, FABL.lowDegreeJuntaCutoff, FABL.lowDegreeJuntaDyadicNodeFailure, FABL.lowDegreeJuntaFourierBatchLinearCoefficient, FABL.lowDegreeJuntaFourierBatchSize_le_linearConfidence, FABL.lowDegreeJuntaFourierFamily, FABL.lowDegreeJuntaNodeAnalysisLinearCoefficient, FABL.lowDegreeJuntaNodeAnalysisWorkBound, FABL.lowDegreeJuntaNodeBaseResourceExponent, FABL.lowDegreeJuntaPolynomialDegree, FABL.lowDegreeJuntaPolynomialWorkBound, FABL.lowDegreeJuntaPolynomialWorkBound_le_linearConfidence, FABL.lowDegreeJuntaPolynomialWorkLinearCoefficient, FABL.lowDegreeJuntaRuntimeExponentLedger}
\difficulty{5}
Let $n,k$, and $b$ be natural numbers, and let $P$ be a set of coordinates among the $n$
coordinates. The node analysis work bound for $P$ at parameter $k$, with failure
parameter equal to the per-call failure obtained from the dyadic failure parameter with
$b$ bits, is at most the node analysis linear coefficient at $k$ times $(n+1)^{\lfloor
3k/4\rfloor+3}(b+1)$.
\end{theorem}

\begin{theorem}[Node Random Example Bound]
\lean{FABL.lowDegreeJuntaNodeRandomExamples_le_linearConfidence}
\label{FABL.lowDegreeJuntaNodeRandomExamples_le_linearConfidence}
\leanok
\uses{FABL.fourierEstimatorFailureBits, FABL.juntaRestrictionSampleCount, FABL.juntaRestrictionSampleCount_cast_le, FABL.lowDegreeJuntaDyadicNodeFailure, FABL.lowDegreeJuntaHalfFailure, FABL.lowDegreeJuntaHalfNodeFailureBits_le, FABL.lowDegreeJuntaLinearConfidenceCoefficient, FABL.lowDegreeJuntaNodeBaseResourceExponent, FABL.lowDegreeJuntaNodeRandomExampleLinearCoefficient, FABL.lowDegreeJuntaNodeRandomExamples, FABL.lowDegreeJuntaRestrictedSampleCount, FABL.lowDegreeJuntaRestrictedSampleCount_le_linearConfidence, FABL.lowDegreeJuntaRestrictedSampleLinearCoefficient}
\difficulty{4}
For every natural number $n$, every subset $P\subseteq[n]$, and all natural numbers
$k,b$, the node's random-example count for $P$, using failure parameter $\delta_{k,b}$,
is at most $A_k(n+1)^{\lfloor 3k/4\rfloor+3}(b+1)$. Define $T_0=1$, $T_{j+1}=1+2T_j$,
$c_k=\lfloor3k/4\rfloor$, $d_k=k-c_k$, $L_k=\lceil\log_2 T_k\rceil+2(c_k+1)+5$,
$B_k=(c_k+1)(36\cdot4^kL_k)+2^{d_k}(d_k+1+L_k)$, and $A_k=2^{k+1}B_k+16\cdot4^kL_k$; the
per-call failure parameter is $\delta_{k,b}=T_k^{-1}2^{-(b+1)}$, and the node's
random-example count uses the restriction sample count with failure parameter
$\delta_{k,b}/2$.
\end{theorem}

\begin{theorem}[Node Work Bound]
\lean{FABL.lowDegreeJuntaNodeWorkBound_le_linearConfidence}
\label{FABL.lowDegreeJuntaNodeWorkBound_le_linearConfidence}
\leanok
\uses{FABL.lowDegreeJuntaCutoff, FABL.lowDegreeJuntaDyadicNodeFailure, FABL.lowDegreeJuntaNodeAnalysisLinearCoefficient, FABL.lowDegreeJuntaNodeAnalysisWorkBound, FABL.lowDegreeJuntaNodeAnalysisWorkBound_le_linearConfidence, FABL.lowDegreeJuntaNodeBaseResourceExponent, FABL.lowDegreeJuntaNodeRandomExampleLinearCoefficient, FABL.lowDegreeJuntaNodeRandomExamples, FABL.lowDegreeJuntaNodeRandomExamples_le_linearConfidence, FABL.lowDegreeJuntaNodeWorkBound, FABL.lowDegreeJuntaNodeWorkLinearCoefficient, FABL.lowDegreeJuntaNodeWorkResourceExponent, FABL.lowDegreeJuntaRestrictedSampleCount, FABL.lowDegreeJuntaRestrictedSampleCount_le_linearConfidence, FABL.lowDegreeJuntaRestrictedSampleLinearCoefficient}
\difficulty{5}
For every natural number $n$, every subset $P\subseteq[n]$, and all natural numbers
$k,b$, the node work bound at failure parameter $\delta_{k,b}$ is at most
$W_k(n+1)^{\lfloor3k/4\rfloor+4}(b+1)$. Here $T_0=1$, $T_{j+1}=1+2T_j$,
$c_k=\lfloor3k/4\rfloor$, $d_k=k-c_k$, $L_k=\lceil\log_2T_k\rceil+2(c_k+1)+5$, and
$\delta_{k,b}=T_k^{-1}2^{-(b+1)}$. The coefficient is $W_k=2A_k+R_k+S_k$, where
$A_k=2^{k+1}R_k+16\cdot4^kL_k$, $R_k=(c_k+1)(36\cdot4^kL_k)+2^{d_k}(d_k+1+L_k)$,
$S_k=(c_k+1)(36\cdot4^kL_k)(c_k+1)+Q_k+(c_k+1)+(d_k+1)$, and $Q_k=2^{d_k}(d_k+1+L_k)$.
\end{theorem}

\begin{theorem}[Uniform Random Example Bound]
\lean{FABL.lowDegreeJuntaNodeRandomExampleBound_le_linearConfidence}
\label{FABL.lowDegreeJuntaNodeRandomExampleBound_le_linearConfidence}
\leanok
\uses{FABL.lowDegreeJuntaDyadicNodeFailure, FABL.lowDegreeJuntaNodeBaseResourceExponent, FABL.lowDegreeJuntaNodeRandomExampleBound, FABL.lowDegreeJuntaNodeRandomExampleLinearCoefficient, FABL.lowDegreeJuntaNodeRandomExamples_le_linearConfidence}
\difficulty{3}
For all natural numbers $n,k,b$, the maximum, over all subsets $P\subseteq[n]$, of the
node's random-example count at failure parameter $\delta_{k,b}$ is at most
$A_k(n+1)^{\lfloor 3k/4\rfloor+3}(b+1)$. Here $T_0=1$, $T_{j+1}=1+2T_j$, $c_k=\lfloor
3k/4\rfloor$, $d_k=k-c_k$, $L_k=\lceil\log_2 T_k\rceil+2(c_k+1)+5$, and
$A_k=2^{k+1}B_k+16\cdot4^kL_k$, where $B_k=(c_k+1)(36\cdot4^kL_k)+2^{d_k}(d_k+1+L_k)$.
The parameter is $\delta_{k,b}=(1/T_k)2^{-(b+1)}$, and the random-example count for a
fixed $P$ is the restriction sample count with failure parameter $\delta_{k,b}/2$.
\end{theorem}

\begin{theorem}[Uniform Node Work Bound]
\lean{FABL.lowDegreeJuntaNodeUniformWorkBound_le_linearConfidence}
\label{FABL.lowDegreeJuntaNodeUniformWorkBound_le_linearConfidence}
\leanok
\uses{FABL.lowDegreeJuntaDyadicNodeFailure, FABL.lowDegreeJuntaNodeUniformWorkBound, FABL.lowDegreeJuntaNodeWorkBound_le_linearConfidence, FABL.lowDegreeJuntaNodeWorkLinearCoefficient, FABL.lowDegreeJuntaNodeWorkResourceExponent}
\difficulty{3}
For all natural numbers $n,k,b$, the maximum, over all subsets $P\subseteq[n]$, of the
node work bound at failure parameter $\delta_{k,b}$ is at most
$W_k(n+1)^{\lfloor3k/4\rfloor+4}(b+1)$. Here $T_0=1$, $T_{j+1}=1+2T_j$,
$c_k=\lfloor3k/4\rfloor$, $d_k=k-c_k$, $L_k=\lceil\log_2T_k\rceil+2(c_k+1)+5$,
$\delta_{k,b}=T_k^{-1}2^{-(b+1)}$, and $W_k=2A_k+B_k+C_k$, with
$A_k=2^{k+1}R_k+16\cdot4^kL_k$, $R_k=(c_k+1)(36\cdot4^kL_k)+2^{d_k}(d_k+1+L_k)$,
$B_k=(c_k+1)(36\cdot4^kL_k)(c_k+1)+Q_k+(c_k+1)+(d_k+1)$, and
$C_k=Q_k(d_k+2)+(d_k+1)Q_k(d_k+3)+(d_k+1)(d_k+2)$, where $Q_k=2^{d_k}(d_k+1+L_k)$.
\end{theorem}

\begin{theorem}[Pointwise example budget bound]
\lean{FABL.lowDegreeJuntaTotalRandomExampleBudget_le_linearConfidence}
\label{FABL.lowDegreeJuntaTotalRandomExampleBudget_le_linearConfidence}
\leanok
\uses{FABL.juntaTreeCallCount, FABL.lowDegreeJuntaDyadicFailure, FABL.lowDegreeJuntaDyadicNodeFailure, FABL.lowDegreeJuntaNodeBaseResourceExponent, FABL.lowDegreeJuntaNodeRandomExampleBound_le_linearConfidence, FABL.lowDegreeJuntaNodeRandomExampleLinearCoefficient, FABL.lowDegreeJuntaTotalRandomExampleBudget}
\difficulty{3}
For all $n,k,b\in\mathbb{N}$, the total random-example budget with dyadic failure
parameter $2^{-(b+1)}$ is at most $T_k A_k(n+1)^{\lfloor 3k/4\rfloor+3}(b+1)$, where
$T_k$ is the junta-tree call count and $A_k$ is the specified natural-number
linear-confidence coefficient.
\end{theorem}

\begin{theorem}[Pointwise work budget bound]
\lean{FABL.lowDegreeJuntaTotalWorkBudget_le_linearConfidence}
\label{FABL.lowDegreeJuntaTotalWorkBudget_le_linearConfidence}
\leanok
\uses{FABL.juntaTreeCallCount, FABL.lowDegreeJuntaDyadicFailure, FABL.lowDegreeJuntaDyadicNodeFailure, FABL.lowDegreeJuntaNodeUniformWorkBound_le_linearConfidence, FABL.lowDegreeJuntaNodeWorkLinearCoefficient, FABL.lowDegreeJuntaNodeWorkResourceExponent, FABL.lowDegreeJuntaTotalWorkBudget}
\difficulty{3}
For all $n,k,b\in\mathbb{N}$, the total work budget with dyadic failure parameter
$2^{-(b+1)}$ is at most $T_k W_k(n+1)^{\lfloor 3k/4\rfloor+4}(b+1)$, where $T_k$ is the
junta-tree call count and $W_k$ is the specified natural-number node-work coefficient.
\end{theorem}

\begin{theorem}[Random-example cost bound]
\lean{FABL.lowDegreeJuntaLearningAlgorithm_randomExampleCost_le}
\label{FABL.lowDegreeJuntaLearningAlgorithm_randomExampleCost_le}
\leanok
\uses{FABL.LearningAccuracy, FABL.juntaTreeCallCount, FABL.juntaTreePerCallFailure, FABL.lowDegreeJuntaDyadicFailure, FABL.lowDegreeJuntaLearningAlgorithm, FABL.lowDegreeJuntaLearningFailure, FABL.lowDegreeJuntaLearningFailure_eq_dyadicFailure, FABL.lowDegreeJuntaNodeBaseResourceExponent, FABL.lowDegreeJuntaNodeRandomExampleBound, FABL.lowDegreeJuntaNodeRandomExampleLinearCoefficient, FABL.lowDegreeJuntaTotalRandomExampleBudget, FABL.lowDegreeJuntaTotalRandomExampleBudget_le_linearConfidence}
\difficulty{4}
Let $n,k\in\mathbb{N}$ and let the accuracy parameter $\varepsilon$ be rational with
$0\leq\varepsilon\leq\tfrac12$. Define $T_0=1$ and $T_{j+1}=1+2T_j$, and set $c=\lfloor
3k/4\rfloor$, $L=\lceil\log_2 T_k\rceil+2(c+1)+5$, and $C=(c+1)36\cdot 4^k L$. The
low-degree junta learning algorithm's random-example cost at accuracy $\varepsilon$ is
at most $4T_k C(n+1)^{c+3}$.
\end{theorem}

\begin{theorem}[Work cost of the junta learner]
\lean{FABL.lowDegreeJuntaLearningAlgorithm_workCost_le}
\label{FABL.lowDegreeJuntaLearningAlgorithm_workCost_le}
\leanok
\uses{FABL.LearningAccuracy, FABL.juntaTreeCallCount, FABL.juntaTreePerCallFailure, FABL.lowDegreeJuntaDyadicFailure, FABL.lowDegreeJuntaLearningAlgorithm, FABL.lowDegreeJuntaLearningFailure, FABL.lowDegreeJuntaLearningFailure_eq_dyadicFailure, FABL.lowDegreeJuntaNodeUniformWorkBound, FABL.lowDegreeJuntaNodeWorkLinearCoefficient, FABL.lowDegreeJuntaNodeWorkResourceExponent, FABL.lowDegreeJuntaTotalWorkBudget, FABL.lowDegreeJuntaTotalWorkBudget_le_linearConfidence}
\difficulty{4}
For natural numbers $n$ and $k$, and an accuracy parameter $\varepsilon$ with $0 \le
\varepsilon \le 1/2$, the work cost of the low-degree junta learning algorithm is at
most $4$ times the number of calls in a depth-$k$ junta tree times the node work linear
coefficient at $k$, multiplied by $(n+1)^{\lfloor 3k/4\rfloor+4}$.
\end{theorem}

\begin{theorem}[Random-example cost of the junta learner]
\lean{FABL.lowDegreeJuntaLearningAlgorithm_randomExampleCost_isBigO}
\label{FABL.lowDegreeJuntaLearningAlgorithm_randomExampleCost_isBigO}
\leanok
\uses{FABL.LearningAccuracy, FABL.juntaTreeCallCount, FABL.lowDegreeJuntaLearningAlgorithm, FABL.lowDegreeJuntaLearningAlgorithm_randomExampleCost_le, FABL.lowDegreeJuntaNodeBaseResourceExponent, FABL.lowDegreeJuntaNodeRandomExampleLinearCoefficient}
\difficulty{3}
For every integer $k \ge 0$ and every accuracy parameter $\varepsilon \in [0,1/2]$, the
number of random labeled examples used by the low-degree junta learning algorithm on the
$n$-dimensional Boolean cube is $O\bigl((n+1)^{\lfloor 3k/4\rfloor+3}\bigr)$ as $n \to
\infty$.
\end{theorem}

\begin{theorem}[Work Cost of the Junta Learning Algorithm]
\lean{FABL.lowDegreeJuntaLearningAlgorithm_workCost_isBigO}
\label{FABL.lowDegreeJuntaLearningAlgorithm_workCost_isBigO}
\leanok
\uses{FABL.LearningAccuracy, FABL.juntaTreeCallCount, FABL.lowDegreeJuntaLearningAlgorithm, FABL.lowDegreeJuntaLearningAlgorithm_workCost_le, FABL.lowDegreeJuntaNodeWorkLinearCoefficient, FABL.lowDegreeJuntaNodeWorkResourceExponent}
\difficulty{3}
For each fixed $k\in\mathbb{N}$ and each learning accuracy $\varepsilon\in[0,1/2]$, the
work cost of the low-degree junta learning algorithm, viewed as a function of the input
dimension $n$, is $O\!\left((n+1)^{\lfloor 3k/4\rfloor+4}\right)$ as $n\to\infty$.
\end{theorem}

\begin{theorem}[Linear-confidence example budget]
\lean{FABL.lowDegreeJuntaTotalRandomExampleBudget_isBigO_linearConfidence}
\label{FABL.lowDegreeJuntaTotalRandomExampleBudget_isBigO_linearConfidence}
\leanok
\uses{FABL.juntaTreeCallCount, FABL.lowDegreeJuntaDyadicFailure, FABL.lowDegreeJuntaNodeBaseResourceExponent, FABL.lowDegreeJuntaNodeRandomExampleLinearCoefficient, FABL.lowDegreeJuntaTotalRandomExampleBudget, FABL.lowDegreeJuntaTotalRandomExampleBudget_le_linearConfidence}
\difficulty{3}
For each fixed $k\in\mathbb{N}$, as $n$ and $b$ tend jointly to infinity, the total
random-example budget for dimension $n$, parameter $k$, and dyadic failure parameter
$2^{-(b+1)}$ is $O\!\left((n+1)^{\lfloor 3k/4\rfloor+3}(b+1)\right)$.
\end{theorem}

\begin{theorem}[Linear-confidence work budget]
\lean{FABL.lowDegreeJuntaTotalWorkBudget_isBigO_linearConfidence}
\label{FABL.lowDegreeJuntaTotalWorkBudget_isBigO_linearConfidence}
\leanok
\uses{FABL.juntaTreeCallCount, FABL.lowDegreeJuntaDyadicFailure, FABL.lowDegreeJuntaNodeWorkLinearCoefficient, FABL.lowDegreeJuntaNodeWorkResourceExponent, FABL.lowDegreeJuntaTotalWorkBudget, FABL.lowDegreeJuntaTotalWorkBudget_le_linearConfidence}
\difficulty{3}
For each fixed $k\in\mathbb{N}$, as $n$ and $b$ tend jointly to infinity, the total work
budget for dimension $n$, parameter $k$, and dyadic failure parameter $2^{-(b+1)}$ is
$O\!\left((n+1)^{\lfloor 3k/4\rfloor+4}(b+1)\right)$.
\end{theorem}

\begin{definition}[Low-degree junta polynomial resource scale]
\lean{FABL.lowDegreeJuntaPolynomialResourceScale}
\label{FABL.lowDegreeJuntaPolynomialResourceScale}
\leanok
For natural numbers $n$, $c$, $d$, and $e$, define the resource scale by
$(n+1)^d(c+1)^e$.
\end{definition}

\begin{definition}[Node polynomial resource witness]
\lean{FABL.LowDegreeJuntaNodePolynomialWitness}
\label{FABL.LowDegreeJuntaNodePolynomialWitness}
\leanok
\uses{FABL.juntaTreePerCallFailure, FABL.lowDegreeJuntaDyadicFailure, FABL.lowDegreeJuntaNodeRandomExampleBound, FABL.lowDegreeJuntaNodeUniformWorkBound, FABL.lowDegreeJuntaPolynomialResourceScale}
For natural numbers $k$, $a$, and $b$, a node polynomial resource witness consists of a
natural-number constant $C$ such that, for every pair of natural numbers $n$ and $c$,
both the node random-example bound at degree parameter $k$ and per-call failure
parameter given by the tree failure transformation applied to the dyadic failure at $c$,
and the node uniform-work bound with those same parameters, are at most
$C(n+1)^a(c+1)^b$.
\end{definition}

\begin{theorem}[Example budget from polynomial bounds]
\lean{FABL.lowDegreeJuntaTotalRandomExampleBudget_le_of_polynomialWitness}
\label{FABL.lowDegreeJuntaTotalRandomExampleBudget_le_of_polynomialWitness}
\leanok
\uses{FABL.LowDegreeJuntaNodePolynomialWitness, FABL.juntaTreeCallCount, FABL.juntaTreePerCallFailure, FABL.lowDegreeJuntaDyadicFailure, FABL.lowDegreeJuntaNodeRandomExampleBound, FABL.lowDegreeJuntaPolynomialResourceScale, FABL.lowDegreeJuntaTotalRandomExampleBudget}
\difficulty{3}
Fix $k,d,e\in\mathbb{N}$ and suppose there is a constant $C\in\mathbb{N}$ such that, for
every $n,b\in\mathbb{N}$, the uniform node random-example bound at per-call failure
parameter $2^{-(b+1)}/T_k$ is at most $C(n+1)^d(b+1)^e$, where $T_k$ is the junta-tree
call count. Then for every $n,b\in\mathbb{N}$, the total random-example budget at
failure parameter $2^{-(b+1)}$ is at most $T_k C(n+1)^d(b+1)^e$.
\end{theorem}

\begin{theorem}[Work budget from polynomial bounds]
\lean{FABL.lowDegreeJuntaTotalWorkBudget_le_of_polynomialWitness}
\label{FABL.lowDegreeJuntaTotalWorkBudget_le_of_polynomialWitness}
\leanok
\uses{FABL.LowDegreeJuntaNodePolynomialWitness, FABL.juntaTreeCallCount, FABL.juntaTreePerCallFailure, FABL.lowDegreeJuntaDyadicFailure, FABL.lowDegreeJuntaNodeUniformWorkBound, FABL.lowDegreeJuntaPolynomialResourceScale, FABL.lowDegreeJuntaTotalWorkBudget}
\difficulty{3}
Fix $k,d,e\in\mathbb{N}$ and suppose there is a constant $C\in\mathbb{N}$ such that, for
every $n,b\in\mathbb{N}$, both the uniform node random-example bound and the uniform
node work bound at per-call failure parameter $2^{-(b+1)}/T_k$ are at most
$C(n+1)^d(b+1)^e$, where $T_k$ is the junta-tree call count. Then, for every
$n,b\in\mathbb{N}$, the total work budget at failure parameter $2^{-(b+1)}$ is at most
$T_k C(n+1)^d(b+1)^e$.
\end{theorem}

\begin{theorem}[Example budget from polynomial bounds]
\lean{FABL.lowDegreeJuntaTotalRandomExampleBudget_isBigO_of_polynomialWitness}
\label{FABL.lowDegreeJuntaTotalRandomExampleBudget_isBigO_of_polynomialWitness}
\leanok
\uses{FABL.LowDegreeJuntaNodePolynomialWitness, FABL.juntaTreeCallCount, FABL.lowDegreeJuntaDyadicFailure, FABL.lowDegreeJuntaPolynomialResourceScale, FABL.lowDegreeJuntaTotalRandomExampleBudget, FABL.lowDegreeJuntaTotalRandomExampleBudget_le_of_polynomialWitness}
\difficulty{3}
Fix $k,d,e\in\mathbb{N}$ and suppose there is a constant $C\in\mathbb{N}$ such that, for
every $n,b\in\mathbb{N}$, the uniform node random-example bound at per-call failure
parameter $2^{-(b+1)}/T_k$ is at most $C(n+1)^d(b+1)^e$, where $T_k$ is the junta-tree
call count. Then, as $n$ and $b$ tend jointly to infinity, the total random-example
budget at failure parameter $2^{-(b+1)}$ is $O\!\left((n+1)^d(b+1)^e\right)$.
\end{theorem}

\begin{theorem}[Work budget from polynomial bounds]
\lean{FABL.lowDegreeJuntaTotalWorkBudget_isBigO_of_polynomialWitness}
\label{FABL.lowDegreeJuntaTotalWorkBudget_isBigO_of_polynomialWitness}
\leanok
\uses{FABL.LowDegreeJuntaNodePolynomialWitness, FABL.juntaTreeCallCount, FABL.lowDegreeJuntaDyadicFailure, FABL.lowDegreeJuntaPolynomialResourceScale, FABL.lowDegreeJuntaTotalWorkBudget, FABL.lowDegreeJuntaTotalWorkBudget_le_of_polynomialWitness}
\difficulty{3}
Fix $k,d,e\in\mathbb{N}$ and suppose there is a constant $C\in\mathbb{N}$ such that, for
every $n,b\in\mathbb{N}$, both the uniform node random-example bound and the uniform
node work bound at per-call failure parameter $2^{-(b+1)}/T_k$ are at most
$C(n+1)^d(b+1)^e$, where $T_k$ is the junta-tree call count. Then, as $n$ and $b$ tend
jointly to infinity, the total work budget at failure parameter $2^{-(b+1)}$ is
$O\!\left((n+1)^d(b+1)^e\right)$.
\end{theorem}
